\SourcePage{575}{578}
\section{Nuclear forces}

The \emph{nuclear forces} specific to nucleons are distinguished above all by their short range: at distances of order $10^{-13}$ cm they fall off exponentially.


In the nonrelativistic limit, nuclear forces may be regarded as independent of the nucleon velocities and as possessing a potential (inside a nucleus, nucleon velocities are about one quarter of the speed of light; see below).
The interaction potential energy $U$ of two nucleons depends not only on their separation $r$, but also, and quite strongly, on their spins\footnote{This sharply distinguishes the interaction of nucleons from that of electrons: electron spin--spin interaction is purely relativistic in origin and is small in atoms.}.
Only a systematic theory of nuclear forces could, of course, determine the exact dependence on $r$.
The form of the spin dependence, however, follows already from simple arguments based on the properties of spin operators.


Only three vectors are available on which the interaction energy $U$ can
depend: the unit vector $\mathbf n$ directed along the radius vector between
the two nucleons, and their spins $\mathbf s_1$ and $\mathbf s_2$. By the
general properties of the spin-$1/2$ operator, every function of it reduces to
a linear function (see §~55). We must also note that the product $\mathbf
n\mathbf s$ is a pseudoscalar rather than a true scalar (since $\mathbf n$ is
a polar vector and $\mathbf s$ an axial one). It is therefore evident that
only two independent scalar quantities, each linear in either spin, can be
formed from $\mathbf n$, $\mathbf s_1$, and $\mathbf s_2$: $\mathbf s_1\mathbf
s_2$ and $(\mathbf n\mathbf s_1)(\mathbf n\mathbf s_2)$\footnote{It is
understood here that the nuclear forces are invariant under spatial inversion,
and hence cannot contain pseudoscalar terms. There is at present no experimental
evidence indicating otherwise.}.

Consequently, as regards its spin dependence, the interaction operator for two
nucleons can be written as a sum of three independent terms,
\begin{equation}
 \hat U_{\text{ord}}=U_1(r)+U_2(r)(\hat{\mathbf s}_1\hat{\mathbf s}_2)
 +U_3(r)\bigl[3(\hat{\mathbf s}_1\mathbf n)
 (\hat{\mathbf s}_2\mathbf n)-\hat{\mathbf s}_1\hat{\mathbf s}_2\bigr],
 \tag{117.1}
\end{equation}
one of which is spin-independent while the other two depend on the spins. The
third term has been written in this form so that it vanishes upon averaging
over the directions of $\mathbf n$; the forces represented by this term are
usually called \emph{tensor forces}.

We have attached the label ``ordinary'' to the interaction (117.1) in order to
emphasize that this operator does not change the

\SourcePage{576}{579}
charge state of the nucleons. An interaction in which a proton turns into a
neutron, and conversely, is also possible. The operator of this ``exchange''
interaction differs in form from (117.1) through the presence of the particle
permutation operator (116.4):
\begin{equation}
 \hat U_{\text{ex}}=\left\{U_4(r)+U_5(r)
 (\hat{\mathbf s}_1\hat{\mathbf s}_2)+U_6(r)
 \bigl[3(\hat{\mathbf s}_1\mathbf n)(\hat{\mathbf s}_2\mathbf n)
 -\hat{\mathbf s}_1\hat{\mathbf s}_2\bigr]\right\}\hat P.
 \tag{117.2}
\end{equation}
The complete interaction operator is the sum
\begin{equation}
 \hat U=\hat U_{\text{ord}}+\hat U_{\text{ex}}.
 \tag{117.3}
\end{equation}
Thus the interaction of two nucleons is specified by six different functions
of their separation. In general, all these terms have the same order of
magnitude\footnote{We also mention that an interaction depending on the
nucleon velocities is described, to first order in the velocities, by an
operator of the form
\[
 [\varphi_1(r)+\varphi_2(r)\hat P]\hat{\mathbf L}\hat{\mathbf S},
\]
where $\mathbf L=\mathbf r\times\mathbf p$ is the orbital angular momentum of
the relative nucleon motion ($\mathbf p$ is its momentum), and $\mathbf
S=\mathbf s_1+\mathbf s_2$; this operator contains two functions of $r$. Terms
of the forms $\mathbf p\mathbf n$ and $\mathbf S\mathbf n$ are excluded by
invariance under spatial inversion and time reversal.}.

The spin operators occurring in (117.1) and (117.2) may be expressed in terms
of the total-spin operator $\mathbf S$. Indeed, on squaring the equalities
$\hat{\mathbf S}=\hat{\mathbf s}_1+\hat{\mathbf s}_2$ and
$\hat{\mathbf S}\mathbf n=\hat{\mathbf s}_1\mathbf n+
\hat{\mathbf s}_2\mathbf n$, and using
$\hat{\mathbf s}_1^2=\hat{\mathbf s}_2^2=3/4$ and
$(\hat{\mathbf s}_1\mathbf n)^2=(\hat{\mathbf s}_2\mathbf n)^2=1/4$
(see (55.10)), we obtain
\begin{equation}
 \hat{\mathbf s}_1\hat{\mathbf s}_2
 =\frac12\left(\hat{\mathbf S}^2-\frac32\right),\qquad
 (\hat{\mathbf s}_1\mathbf n)(\hat{\mathbf s}_2\mathbf n)
 =\frac12\left[(\hat{\mathbf S}\mathbf n)^2-\frac12\right].
 \tag{117.4}
\end{equation}

The operator $\hat{\mathbf S}^2$ commutes with $\hat{\mathbf S}$; hence the
interactions represented by the first two terms of (117.1) and (117.2)
preserve the system's total-spin vector. The tensor interaction, by contrast,
contains the operator $(\hat{\mathbf S}\mathbf n)^2$, which commutes with the
square $\hat{\mathbf S}^2$ but not with the vector $\hat{\mathbf S}$ itself.
It follows that only the magnitude of the total spin remains conserved, not
its direction.

For a system of two nucleons the total spin $S$ may take the values 0 and 1.
Its total isotopic spin $T$ has the same two possible values. All states of the
system therefore separate into four groups, distinguished by the pair of
numbers $S,T$. For the states in each group

\SourcePage{577}{580}
For $S=0$, the corresponding interaction operator has the form $A(r)$.
For $S=1$, it has the form
$A(r)+B(r)[(\hat{\mathbf S}\mathbf n)^2-2/3]$.
In these cases the general operator (117.3) reduces to the respective one
(see Problem 1)\footnote{The sign of the observed quadrupole moment of the deuteron implies that, in the $T=0$, $S=1$ state, the coefficient $B(r)$ of the tensor forces is negative.
Experimental data on the properties of the deuteron further show that the nucleon interaction in this state contains a strong attraction with a deep potential well.
The presence of tensor forces makes it difficult to express this fact in terms of properties of the functions $A(r)$ and $B(r)$.
Nucleon-scattering data show that attraction is also present for $T=1$, $S=0$.
It is weaker and, in particular, does not result in a stable two-particle system.
}.


For fixed $S$ and $T$, the states of the system are classified by total angular momentum $J$ and parity.
As we know, $T=0$ corresponds to states with symmetric wave functions $\psi$, whereas for $T=1$ these functions are antisymmetric.
On the other hand, $S$ fixes the symmetry of the wave function in the spin variables: it is symmetric for $S=1$ and antisymmetric for $S=0$.
Consequently, the pair $S,T$ also fixes the symmetry in the spatial variables, and hence the parity of the state.
Thus, a system with isotopic spin $T=0$ can have only even triplet states ($S=1$) or odd singlet states ($S=0$); for $T=1$, the states are odd triplets or even singlets.

Since the spin vector is not conserved, orbital angular momentum is not generally conserved either; only the sum $\mathbf J=\mathbf L+\mathbf S$ is conserved.
Its magnitude $L$ may nevertheless be conserved simply because fixed values of $J$, $S$, and parity (or of $J$, $S$, and $T$) can be compatible with just one definite value of $L$; recall that the parity of a two-particle system is $(-1)^L$.

Thus an odd state with $S=1$, $J=1$
can have only $L=1$, and is therefore a ${}^3P_1$ state. In other cases, two
different values of $L$ may correspond to the prescribed values of $J$, $S$,
and parity, so that $L$ is not conserved. For example, in an odd state with
$S=1$, $J=2$, either $L=1$ or $L=3$ is possible; such a state is a
superposition ${}^3P_2+{}^3F_2$.

We thus arrive at the following possible states of a two-nucleon system (the
sign $\pm$ indicates the parity):

\SourcePage{578}{581}
\begin{align*}
 \text{for }T=1:&\quad {}^3P_0^-,\ {}^3P_1^-,\
 ({}^3P_2+{}^3F_2)^-,\ {}^3F_3^-,\ldots;\\
 &\quad {}^1S_0^+,\ {}^1D_2^+,\ {}^1G_4^+,\ldots,\\
 \text{for }T=0:&\quad ({}^3S_1+{}^3D_1)^+,\ {}^3D_2^+,\\
 &\quad ({}^3D_3+{}^3G_3)^+,\ldots;\
 {}^1P_1^-,\ {}^1F_3^-,\ldots
\end{align*}

In general, nuclear forces are nonadditive. This means that the interaction in
a system containing more than two nucleons does not reduce to the sum of the
interactions between all particle pairs. It appears, however, that three-body
and higher interactions play a relatively minor role compared with pair
interactions. The properties of complex nuclei can therefore be treated to a
considerable extent on the basis of pairwise interactions.

Experimental information on nuclei shows that, as the particle number $A$
increases, a nucleon system begins to behave as macroscopic ``nuclear matter,''
whose volume and energy grow in proportion to $A$ (apart from effects due to
the Coulomb interaction of the protons and to the presence of a free nuclear
surface). The property of nuclear forces responsible for this behavior is
called their \emph{saturation}.

This saturation property imposes certain restrictions on the functions $U_1,\ldots,U_6$ that specify pairwise nucleon interactions.
Suppose all the particles are confined within a region whose dimensions are of the order of the range of the nuclear forces; every pair then interacts.
If some configuration of the nucleons, including their spin orientations, makes all these pair forces attractive, the potential energy of the system is negative and proportional to $A^2$.
Its kinetic energy, by contrast, is positive and proportional to $A^{5/3}$, a lower power of $A$\footnote{The density $n$ of particles confined in the given volume is proportional to their number $A$, while the kinetic energy per particle is then proportional to $n^{2/3}$ (compare (70.1)). The total kinetic energy is therefore proportional to $AA^{2/3}$.}.
Under these conditions, a sufficiently large collection of nucleons would indeed contract into a small volume independent of $A$, and hence would not form nuclear matter.
Saturation of the nuclear forces must therefore require the absence of configurations whose negative interaction energy is proportional to $A^2$ (see problem 2).

The fact that the volume of nuclear matter is proportional to the particle number is expressed by a relation of the form

\begin{equation}
 R=r_0A^{1/3},
 \tag{117.5}
\end{equation}

\SourcePage{579}{582}
which connects the nuclear radius $R$ with the number $A$ of particles in the
nucleus. Experimental data (from electron scattering by nuclei) give
$r_0\approx1{.}1\cdot10^{-13}$ cm.

Let us determine the limiting nucleon momentum in nuclear matter (compare
§~70). The phase-space volume corresponding to particles that occupy a unit
volume of physical space and have momenta $p\leqslant p_0$ is $4\pi p_0^3/3$.
Dividing this by $(2\pi\hbar)^3$ gives the number of ``cells,'' each capable of
holding two protons and two neutrons. Taking the numbers of protons and
neutrons to be equal, we obtain
\[
 4\frac{4\pi}{3}\left(\frac{p_0}{2\pi\hbar}\right)^3=\frac AV,
\]
where $V$ is the nuclear volume. Substitution of (117.5) then gives
\[
 p_0=\left(\frac{3\pi^2A}{2V}\right)^{1/3}\hbar
 =\frac{(9\pi)^{1/3}\hbar}{2r_0}
 =1{.}4\cdot10^{-14}\ \text{g}\cdot\text{cm}\cdot\text{s}^{-1}.
\]
The corresponding energy $p_0^2/2m_p$ (where $m_p$ is the nucleon mass) is
$\sim30$ MeV, and the velocity is $p_0/m_p\approx c/4$.

\subsection*{Problems}

\ProblemHeading{1.} Find the interaction operators for two nucleons in states
with definite values of $S$ and $T$.

\SolutionHeading The required operators $\hat U_{ST}$ follow from the general
expressions (117.1)--(117.3), with relations (116.3) and (117.4) taken into
account:
\begin{align*}
 \hat U_{00}&=U_1-\frac34U_2+U_4-\frac34U_5,\\
 \hat U_{01}&=U_1-\frac34U_2-U_4+\frac34U_5,\\
 \hat U_{10}&=U_1+\frac14U_2+U_4+\frac14U_5
 +\frac12(U_3+U_6)\bigl[3(\hat{\mathbf S}\mathbf n)^2-2\bigr],\\
 \hat U_{11}&=U_1+\frac14U_2-U_4-\frac14U_5
 +\frac12(U_3-U_6)\bigl[3(\hat{\mathbf S}\mathbf n)^2-2\bigr].
\end{align*}

\ProblemHeading{2.} Find the saturation conditions for the nuclear forces,
assuming that tensor forces are absent and that the ranges of all the remaining
types of force are equal.

\SolutionHeading Consider several limiting cases (between which all other
possible cases lie) for the state of a system of $A$ nucleons, and write the
conditions under which the interaction energy of an ``average'' nucleon pair
in this system is positive.

Let the total spin and isotopic spin of the nucleus have their largest possible
values: $S_{\text{nuc}}=T_{\text{nuc}}=A/2$ (all particles in the system are
protons with parallel spins). Every nucleon pair then has $S=T=1$, giving the
condition
\begin{equation}
 U_{11}>0.
 \tag{1}
\end{equation}

\SourcePage{580}{583}
Now let $T_{\text{nuc}}=A/2$, $S_{\text{nuc}}=0$. Every nucleon pair then has
$T=1$, while the mean value of $s_z$ for a single nucleon is zero. The latter
means that a nucleon is equally likely to have $s_z=1/2$ or $s_z=-1/2$; under
these conditions the probabilities for a nucleon pair to be in states with
$S=0$ and $S=1$ are respectively $1/2$ and $3/4$ (they are proportional to
the number $2S+1$ of possible $S_z$ values). The condition that the mean pair
energy be positive is therefore
\begin{equation}
 \frac14U_{01}+\frac34U_{11}>0.
 \tag{2}
\end{equation}

Similarly, consideration of the state with $T_{\text{nuc}}=0$,
$S_{\text{nuc}}=A/2$ leads to
\begin{equation}
 \frac14U_{10}+\frac34U_{11}>0.
 \tag{3}
\end{equation}
In the state with $T_{\text{nuc}}=S_{\text{nuc}}=0$, the probability that a
nucleon pair has $S=T=1$ is $3/4\cdot3/4$, the probability that it has $T=1$,
$S=0$ is $3/4\cdot1/4$, and so on. We thus find the condition
\begin{equation}
 \frac9{16}U_{11}+\frac3{16}(U_{10}+U_{01})+\frac1{16}U_{00}>0.
 \tag{4}
\end{equation}

Finally, let the system contain $A/2$ protons and $A/2$ neutrons, with the
spins of all protons parallel to one another and antiparallel to the spins of
all neutrons. A given nucleon is equally likely to be a $p$ or an $n$, that is,
to have $T_\zeta=1/2$ or $T_\zeta=-1/2$; the probability that a nucleon pair
has $T=0$ is $1/4$. In this case one nucleon in the pair is a $p$ and the
other an $n$, and hence $S_z=0$. This value of $S_z$ may arise with equal
probability from states having $S=0$ or $S=1$. Consequently, the probabilities
for the pair to occupy a state with $T=0$, $S=0$ or with $T=0$, $S=1$ are
each $1/4\cdot1/2=1/8$. The probability for the state $T=1$, $S=0$ is the
same, while the remaining $5/8$ belongs to the state $T=S=1$. Taking all this
into account, we obtain
\begin{equation}
 \frac18(U_{00}+U_{01}+U_{10})+\frac58U_{11}>0.
 \tag{5}
\end{equation}
Inequalities (1)--(5) constitute the required system of saturation conditions
for the nuclear forces.
